\documentclass[11pt]{article}
% ============================================================================
% preamble.tex -- shared style for every document in docs/
%
% Each document starts with:
%     \documentclass[11pt]{article}
%     % ============================================================================
% preamble.tex -- shared style for every document in docs/
%
% Each document starts with:
%     \documentclass[11pt]{article}
%     % ============================================================================
% preamble.tex -- shared style for every document in docs/
%
% Each document starts with:
%     \documentclass[11pt]{article}
%     \input{preamble}
%     \newcommand{\doctitle}{...}
%
% Team names live in ONE place (below). Edit them here and every document
% updates on the next build.
% ============================================================================

% ---------------------------------------------------------------- team
\newcommand{\projname}{FPGA int8 MLP Accelerator}
\newcommand{\repourl}{https://github.com/itsKannapi/fpga-mlp-accelerator}
\newcommand{\memberA}{Joey (\texttt{itsKannapi})}
\newcommand{\memberB}{Teammate~2 (TBD)}
\newcommand{\memberC}{Teammate~3 (TBD)}

% ---------------------------------------------------------------- packages
\usepackage[margin=1in]{geometry}
\usepackage[T1]{fontenc}
\usepackage{lmodern}
\usepackage{microtype}
\usepackage{amsmath}
\usepackage{booktabs}
\usepackage{tabularx}
\usepackage{longtable}
\usepackage{array}
\usepackage{enumitem}
\usepackage{xcolor}
\usepackage{listings}
\usepackage{tikz}
\usetikzlibrary{arrows.meta, positioning, fit, calc}
\usepackage[most]{tcolorbox}
\usepackage{fancyhdr}
\usepackage{titlesec}
\usepackage{needspace}
\usepackage[hidelinks]{hyperref}
\hypersetup{colorlinks=true, linkcolor=accent, urlcolor=accent, citecolor=accent}

% ---------------------------------------------------------------- colors
\definecolor{accent}{HTML}{1F4E79}
\definecolor{muted}{HTML}{5F6B7A}
\definecolor{rule}{HTML}{C9D1DA}
\definecolor{okc}{HTML}{2E7D32}
\definecolor{progc}{HTML}{B26A00}
\definecolor{todoc}{HTML}{8A8F98}
\definecolor{codebg}{HTML}{F5F7FA}

% ---------------------------------------------------------------- layout
\setcounter{tocdepth}{1}   % sections only, keeps the contents to a few lines
\setlength{\parindent}{0pt}
\setlength{\parskip}{0.55em}
\setlist{itemsep=0.15em, topsep=0.3em}
\renewcommand{\arraystretch}{1.18}

\titleformat{\section}{\Large\bfseries\color{accent}}{\thesection}{0.7em}{}
\titleformat{\subsection}{\large\bfseries}{\thesubsection}{0.6em}{}
\titleformat{\subsubsection}{\normalsize\bfseries}{\thesubsubsection}{0.5em}{}
\titlespacing*{\section}{0pt}{1.4em}{0.5em}

\pagestyle{fancy}
\fancyhf{}
\fancyhead[L]{\small\color{muted}\projname}
\fancyhead[R]{\small\color{muted}\doctitle}
\fancyfoot[C]{\small\color{muted}\thepage}
\renewcommand{\headrulewidth}{0.4pt}
\renewcommand{\headrule}{\color{rule}\hrule width\headwidth height\headrulewidth}

% ---------------------------------------------------------------- title block
\newcommand{\makedoctitle}[2]{%  #1 = subtitle, #2 = last-updated date
  \thispagestyle{plain}
  {\color{muted}\small\MakeUppercase{\projname}}\par\vspace{0.2em}
  {\Huge\bfseries\color{accent}\doctitle}\par\vspace{0.4em}
  {\large #1}\par\vspace{0.6em}
  {\small\color{muted}Team: \memberA, \memberB, \memberC\quad$\cdot$\quad
   Last updated #2\quad$\cdot$\quad\href{\repourl}{\texttt{\repourl}}}\par
  \vspace{0.4em}{\color{rule}\rule{\linewidth}{0.6pt}}\par\vspace{0.6em}
}

% ---------------------------------------------------------------- status tags
\newcommand{\tagbox}[2]{\tcbox[on line, colback=#1!12, colframe=#1!12,
  coltext=#1!70!black, boxsep=0pt, left=3pt, right=3pt, top=1.5pt, bottom=1.5pt,
  arc=2pt]{\footnotesize\bfseries #2}}
\newcommand{\done}{\tagbox{okc}{DONE}}
\newcommand{\inprog}{\tagbox{progc}{IN PROGRESS}}
\newcommand{\todo}{\tagbox{todoc}{TO DO}}
\newcommand{\stretchgoal}{\tagbox{accent}{STRETCH}}

% ---------------------------------------------------------------- callouts
\newtcolorbox{note}[1][Note]{colback=accent!5, colframe=accent!60,
  boxrule=0.6pt, arc=2pt, left=6pt, right=6pt, top=4pt, bottom=4pt,
  fonttitle=\bfseries\small, title=#1, coltitle=white}
\newtcolorbox{warn}[1][Watch out]{colback=progc!6, colframe=progc!70,
  boxrule=0.6pt, arc=2pt, left=6pt, right=6pt, top=4pt, bottom=4pt,
  fonttitle=\bfseries\small, title=#1, coltitle=white}

% ---------------------------------------------------------------- code
\lstdefinelanguage{SV}{
  morekeywords={module,endmodule,input,output,wire,reg,signed,always,posedge,
    negedge,if,else,begin,end,assign,parameter,localparam,generate,endgenerate,
    for,genvar,case,endcase,function,endfunction,integer,initial},
  sensitive=true, morecomment=[l]{//}, morecomment=[s]{/*}{*/},
  morestring=[b]", morestring=[b]`}
\lstset{
  basicstyle=\ttfamily\small, backgroundcolor=\color{codebg},
  keywordstyle=\color{accent}\bfseries, commentstyle=\color{muted}\itshape,
  frame=single, rulecolor=\color{rule}, framesep=4pt, xleftmargin=4pt,
  xrightmargin=4pt, columns=fullflexible, keepspaces=true, breaklines=true,
  showstringspaces=false, upquote=true, aboveskip=0.6em, belowskip=0.6em}
\newcommand{\cmd}[1]{\texttt{#1}}
\newcommand{\file}[1]{\texttt{#1}}

%     \newcommand{\doctitle}{...}
%
% Team names live in ONE place (below). Edit them here and every document
% updates on the next build.
% ============================================================================

% ---------------------------------------------------------------- team
\newcommand{\projname}{FPGA int8 MLP Accelerator}
\newcommand{\repourl}{https://github.com/itsKannapi/fpga-mlp-accelerator}
\newcommand{\memberA}{Joey (\texttt{itsKannapi})}
\newcommand{\memberB}{Teammate~2 (TBD)}
\newcommand{\memberC}{Teammate~3 (TBD)}

% ---------------------------------------------------------------- packages
\usepackage[margin=1in]{geometry}
\usepackage[T1]{fontenc}
\usepackage{lmodern}
\usepackage{microtype}
\usepackage{amsmath}
\usepackage{booktabs}
\usepackage{tabularx}
\usepackage{longtable}
\usepackage{array}
\usepackage{enumitem}
\usepackage{xcolor}
\usepackage{listings}
\usepackage{tikz}
\usetikzlibrary{arrows.meta, positioning, fit, calc}
\usepackage[most]{tcolorbox}
\usepackage{fancyhdr}
\usepackage{titlesec}
\usepackage{needspace}
\usepackage[hidelinks]{hyperref}
\hypersetup{colorlinks=true, linkcolor=accent, urlcolor=accent, citecolor=accent}

% ---------------------------------------------------------------- colors
\definecolor{accent}{HTML}{1F4E79}
\definecolor{muted}{HTML}{5F6B7A}
\definecolor{rule}{HTML}{C9D1DA}
\definecolor{okc}{HTML}{2E7D32}
\definecolor{progc}{HTML}{B26A00}
\definecolor{todoc}{HTML}{8A8F98}
\definecolor{codebg}{HTML}{F5F7FA}

% ---------------------------------------------------------------- layout
\setcounter{tocdepth}{1}   % sections only, keeps the contents to a few lines
\setlength{\parindent}{0pt}
\setlength{\parskip}{0.55em}
\setlist{itemsep=0.15em, topsep=0.3em}
\renewcommand{\arraystretch}{1.18}

\titleformat{\section}{\Large\bfseries\color{accent}}{\thesection}{0.7em}{}
\titleformat{\subsection}{\large\bfseries}{\thesubsection}{0.6em}{}
\titleformat{\subsubsection}{\normalsize\bfseries}{\thesubsubsection}{0.5em}{}
\titlespacing*{\section}{0pt}{1.4em}{0.5em}

\pagestyle{fancy}
\fancyhf{}
\fancyhead[L]{\small\color{muted}\projname}
\fancyhead[R]{\small\color{muted}\doctitle}
\fancyfoot[C]{\small\color{muted}\thepage}
\renewcommand{\headrulewidth}{0.4pt}
\renewcommand{\headrule}{\color{rule}\hrule width\headwidth height\headrulewidth}

% ---------------------------------------------------------------- title block
\newcommand{\makedoctitle}[2]{%  #1 = subtitle, #2 = last-updated date
  \thispagestyle{plain}
  {\color{muted}\small\MakeUppercase{\projname}}\par\vspace{0.2em}
  {\Huge\bfseries\color{accent}\doctitle}\par\vspace{0.4em}
  {\large #1}\par\vspace{0.6em}
  {\small\color{muted}Team: \memberA, \memberB, \memberC\quad$\cdot$\quad
   Last updated #2\quad$\cdot$\quad\href{\repourl}{\texttt{\repourl}}}\par
  \vspace{0.4em}{\color{rule}\rule{\linewidth}{0.6pt}}\par\vspace{0.6em}
}

% ---------------------------------------------------------------- status tags
\newcommand{\tagbox}[2]{\tcbox[on line, colback=#1!12, colframe=#1!12,
  coltext=#1!70!black, boxsep=0pt, left=3pt, right=3pt, top=1.5pt, bottom=1.5pt,
  arc=2pt]{\footnotesize\bfseries #2}}
\newcommand{\done}{\tagbox{okc}{DONE}}
\newcommand{\inprog}{\tagbox{progc}{IN PROGRESS}}
\newcommand{\todo}{\tagbox{todoc}{TO DO}}
\newcommand{\stretchgoal}{\tagbox{accent}{STRETCH}}

% ---------------------------------------------------------------- callouts
\newtcolorbox{note}[1][Note]{colback=accent!5, colframe=accent!60,
  boxrule=0.6pt, arc=2pt, left=6pt, right=6pt, top=4pt, bottom=4pt,
  fonttitle=\bfseries\small, title=#1, coltitle=white}
\newtcolorbox{warn}[1][Watch out]{colback=progc!6, colframe=progc!70,
  boxrule=0.6pt, arc=2pt, left=6pt, right=6pt, top=4pt, bottom=4pt,
  fonttitle=\bfseries\small, title=#1, coltitle=white}

% ---------------------------------------------------------------- code
\lstdefinelanguage{SV}{
  morekeywords={module,endmodule,input,output,wire,reg,signed,always,posedge,
    negedge,if,else,begin,end,assign,parameter,localparam,generate,endgenerate,
    for,genvar,case,endcase,function,endfunction,integer,initial},
  sensitive=true, morecomment=[l]{//}, morecomment=[s]{/*}{*/},
  morestring=[b]", morestring=[b]`}
\lstset{
  basicstyle=\ttfamily\small, backgroundcolor=\color{codebg},
  keywordstyle=\color{accent}\bfseries, commentstyle=\color{muted}\itshape,
  frame=single, rulecolor=\color{rule}, framesep=4pt, xleftmargin=4pt,
  xrightmargin=4pt, columns=fullflexible, keepspaces=true, breaklines=true,
  showstringspaces=false, upquote=true, aboveskip=0.6em, belowskip=0.6em}
\newcommand{\cmd}[1]{\texttt{#1}}
\newcommand{\file}[1]{\texttt{#1}}

%     \newcommand{\doctitle}{...}
%
% Team names live in ONE place (below). Edit them here and every document
% updates on the next build.
% ============================================================================

% ---------------------------------------------------------------- team
\newcommand{\projname}{FPGA int8 MLP Accelerator}
\newcommand{\repourl}{https://github.com/itsKannapi/fpga-mlp-accelerator}
\newcommand{\memberA}{Joey (\texttt{itsKannapi})}
\newcommand{\memberB}{Teammate~2 (TBD)}
\newcommand{\memberC}{Teammate~3 (TBD)}

% ---------------------------------------------------------------- packages
\usepackage[margin=1in]{geometry}
\usepackage[T1]{fontenc}
\usepackage{lmodern}
\usepackage{microtype}
\usepackage{amsmath}
\usepackage{booktabs}
\usepackage{tabularx}
\usepackage{longtable}
\usepackage{array}
\usepackage{enumitem}
\usepackage{xcolor}
\usepackage{listings}
\usepackage{tikz}
\usetikzlibrary{arrows.meta, positioning, fit, calc}
\usepackage[most]{tcolorbox}
\usepackage{fancyhdr}
\usepackage{titlesec}
\usepackage{needspace}
\usepackage[hidelinks]{hyperref}
\hypersetup{colorlinks=true, linkcolor=accent, urlcolor=accent, citecolor=accent}

% ---------------------------------------------------------------- colors
\definecolor{accent}{HTML}{1F4E79}
\definecolor{muted}{HTML}{5F6B7A}
\definecolor{rule}{HTML}{C9D1DA}
\definecolor{okc}{HTML}{2E7D32}
\definecolor{progc}{HTML}{B26A00}
\definecolor{todoc}{HTML}{8A8F98}
\definecolor{codebg}{HTML}{F5F7FA}

% ---------------------------------------------------------------- layout
\setcounter{tocdepth}{1}   % sections only, keeps the contents to a few lines
\setlength{\parindent}{0pt}
\setlength{\parskip}{0.55em}
\setlist{itemsep=0.15em, topsep=0.3em}
\renewcommand{\arraystretch}{1.18}

\titleformat{\section}{\Large\bfseries\color{accent}}{\thesection}{0.7em}{}
\titleformat{\subsection}{\large\bfseries}{\thesubsection}{0.6em}{}
\titleformat{\subsubsection}{\normalsize\bfseries}{\thesubsubsection}{0.5em}{}
\titlespacing*{\section}{0pt}{1.4em}{0.5em}

\pagestyle{fancy}
\fancyhf{}
\fancyhead[L]{\small\color{muted}\projname}
\fancyhead[R]{\small\color{muted}\doctitle}
\fancyfoot[C]{\small\color{muted}\thepage}
\renewcommand{\headrulewidth}{0.4pt}
\renewcommand{\headrule}{\color{rule}\hrule width\headwidth height\headrulewidth}

% ---------------------------------------------------------------- title block
\newcommand{\makedoctitle}[2]{%  #1 = subtitle, #2 = last-updated date
  \thispagestyle{plain}
  {\color{muted}\small\MakeUppercase{\projname}}\par\vspace{0.2em}
  {\Huge\bfseries\color{accent}\doctitle}\par\vspace{0.4em}
  {\large #1}\par\vspace{0.6em}
  {\small\color{muted}Team: \memberA, \memberB, \memberC\quad$\cdot$\quad
   Last updated #2\quad$\cdot$\quad\href{\repourl}{\texttt{\repourl}}}\par
  \vspace{0.4em}{\color{rule}\rule{\linewidth}{0.6pt}}\par\vspace{0.6em}
}

% ---------------------------------------------------------------- status tags
\newcommand{\tagbox}[2]{\tcbox[on line, colback=#1!12, colframe=#1!12,
  coltext=#1!70!black, boxsep=0pt, left=3pt, right=3pt, top=1.5pt, bottom=1.5pt,
  arc=2pt]{\footnotesize\bfseries #2}}
\newcommand{\done}{\tagbox{okc}{DONE}}
\newcommand{\inprog}{\tagbox{progc}{IN PROGRESS}}
\newcommand{\todo}{\tagbox{todoc}{TO DO}}
\newcommand{\stretchgoal}{\tagbox{accent}{STRETCH}}

% ---------------------------------------------------------------- callouts
\newtcolorbox{note}[1][Note]{colback=accent!5, colframe=accent!60,
  boxrule=0.6pt, arc=2pt, left=6pt, right=6pt, top=4pt, bottom=4pt,
  fonttitle=\bfseries\small, title=#1, coltitle=white}
\newtcolorbox{warn}[1][Watch out]{colback=progc!6, colframe=progc!70,
  boxrule=0.6pt, arc=2pt, left=6pt, right=6pt, top=4pt, bottom=4pt,
  fonttitle=\bfseries\small, title=#1, coltitle=white}

% ---------------------------------------------------------------- code
\lstdefinelanguage{SV}{
  morekeywords={module,endmodule,input,output,wire,reg,signed,always,posedge,
    negedge,if,else,begin,end,assign,parameter,localparam,generate,endgenerate,
    for,genvar,case,endcase,function,endfunction,integer,initial},
  sensitive=true, morecomment=[l]{//}, morecomment=[s]{/*}{*/},
  morestring=[b]", morestring=[b]`}
\lstset{
  basicstyle=\ttfamily\small, backgroundcolor=\color{codebg},
  keywordstyle=\color{accent}\bfseries, commentstyle=\color{muted}\itshape,
  frame=single, rulecolor=\color{rule}, framesep=4pt, xleftmargin=4pt,
  xrightmargin=4pt, columns=fullflexible, keepspaces=true, breaklines=true,
  showstringspaces=false, upquote=true, aboveskip=0.6em, belowskip=0.6em}
\newcommand{\cmd}[1]{\texttt{#1}}
\newcommand{\file}[1]{\texttt{#1}}

\newcommand{\doctitle}{Progress Log}

\begin{document}
\makedoctitle{Milestones, dated log, open issues, and lessons learned}{2026-10-16}

\begin{note}[How to update this document]
Update it in the \emph{same pull request} as the work it describes: flip the
milestone status, add a dated entry to the top of the log (newest first), and
rebuild the PDF (\cmd{latexmk -pdf progress.tex} in \file{docs/}). Copy the
entry template from Section~\ref{sec:template}. Keep entries factual: what
changed, how it was verified, the commit.
\end{note}

% ============================================================================
\section{Team}
\begin{center}
\begin{tabularx}{\linewidth}{@{}lX@{}}
\toprule
Member & Focus \\
\midrule
\memberA & Project lead; RTL and verification so far \\
\memberB & TBD \\
\memberC & TBD \\
\bottomrule
\end{tabularx}
\end{center}
Names are set once in \file{docs/preamble.tex} and appear in every document.

% ============================================================================
\section{Milestones}
Target: a working, documented on-board demo within four weeks of starting the RTL.

\begin{center}
\begin{longtable}{@{}p{0.6cm}p{8.4cm}p{1.6cm}p{3.0cm}@{}}
\toprule
Wk & Milestone & Owner & Status \\
\midrule
\endhead
1 & Toolchain: Vivado 2026.1 + license, board smoke test & Joey & \done \\
1 & Train and quantize the MLP; export weights and golden vectors & Joey & \done \\
1 & MAC unit + three-tier self-checking testbench & Joey & \done \\
\midrule
2 & Layer 1: 32 parallel MACs, ReLU + requantize, bit-exact vs.\ golden & Joey & \done \\
2 & Repo made portable (project script, Icarus runner); docs; GitHub & Joey & \inprog \\
2 & Layer 2 (32$\rightarrow$10) + argmax & TBD & \todo \\
2 & Full-network testbench: predicted digit + all 10 scores vs.\ golden & TBD & \todo \\
\midrule
3 & Image buffer (784\,B BRAM) + UART receiver & TBD & \todo \\
3 & UART transmitter for the prediction; define the host protocol & TBD & \todo \\
3 & Top-level module + Basys~3 constraints (\file{.xdc}) & TBD & \todo \\
3 & Timing closure at 100\,MHz & TBD & \todo \\
3 & Host script: send 100 test images, report on-board accuracy & TBD & \todo \\
\midrule
4 & Record Vivado utilization + timing in \file{design.tex} & TBD & \todo \\
4 & README polish, block diagram, demo video & TBD & \todo \\
\midrule
-- & Binary/ternary (XNOR + popcount) variant reusing the pipeline & TBD & \stretchgoal \\
\bottomrule
\end{longtable}
\end{center}

% ============================================================================
\section{Log}
Newest first. Commit hashes refer to the \texttt{main} branch.

\subsection*{2026-10-16 \quad Layer 1 verified; repo prepared for the team}
\begin{itemize}
  \item Built \file{src/layer1.v}: 32 MACs fed by one broadcast pixel per clock and a
        256-bit weight word. Latency 786 cycles (7.86\,\textmu s at 100\,MHz).
  \item Exporter now also writes \file{rtl/w1\_t.mem} (pixel-major). Re-running it
        reproduced every existing file byte for byte.
  \item \file{sim/tb\_layer1.v}: 4{,}259 checks, \textbf{all pass in Vivado xsim}.
        Mutation testing found that real MNIST images miss three bug classes
        (border pixels are always 0, no neuron saturates); added synthetic images, and
        all 10 planted bugs are now caught. Details in \file{design.tex} \S4.
  \item Estimated resources (Yosys): 32 DSP, 15 RAMB18, $\sim$1.7k LUT.
  \item Removed the hard-coded \file{C:/dev/...} path from the Verilog.
  \item Added two scripts: \file{create\_project.tcl} rebuilds the Vivado project on
        any machine, and \file{run\_sims.sh} runs every testbench in Icarus.
        Started these docs.
  \item Commit \texttt{e044fd4} (layer 1); docs/portability commit to follow.
\end{itemize}

\subsection*{2026-10-09 \quad MAC unit verified}
\begin{itemize}
  \item \file{src/mac.v}: int8$\times$int8 into a 32-bit accumulator, mapped to one DSP48.
  \item \file{sim/tb\_mac.v}: 10 directed + 2{,}000 random + 256 golden checks
        (2{,}266 total), \textbf{all pass in Vivado xsim} using the real trained weights.
  \item Three planted bugs each caught; each test tier catches bugs the others miss.
  \item Commit \texttt{1276047}.
\end{itemize}

\subsection*{2026-10-08 \quad Board bring-up; model trained and quantized}
\begin{itemize}
  \item Smoke test: a two-switch AND gate on LD0 programmed over JTAG and working.
        This confirms license, synthesis, implementation, bitstream, drivers.
  \item Trained the $784\rightarrow32\rightarrow10$ MLP: 95.86\% test accuracy.
        int8 version also 95.86\% (99.74\% of predictions identical). \texttt{SHIFT1}~=~12.
  \item Golden vectors for 64 images exported to \file{rtl/golden.txt}.
  \item Commit \texttt{7e70d2d}.
\end{itemize}

\subsection*{2026-10-01 \quad Vivado installed}
\begin{itemize}
  \item Vivado 2026.1 needs a license file even for the free tier. Installed the
        free \emph{Basic Tier, node-locked} license (see \file{setup.tex}).
  \item The Artix-7 family was missing from the install; added it through the
        installer's add-devices mode.
\end{itemize}

\subsection*{Before RTL \quad Project scoped}
Chose an int8 MLP accelerator over a RISC-V CPU and over camera/microphone
projects: it tells a complete ``AI hardware'' story and needs no exotic I/O, so
it is finishable before internship application deadlines.

% ============================================================================
\section{Open issues}
\begin{center}
\begin{tabularx}{\linewidth}{@{}p{0.5cm}Xp{2.2cm}@{}}
\toprule
\# & Issue & Owner \\
\midrule
1 & \file{mlp\_fp32.npz} (37\,MB, holds the float weights) is not in git, so only
    Joey's machine can re-run \file{quantize\_export.py}. Fix before anyone needs
    a new export: commit a small weights-only file and have the script load test
    data itself. & TBD \\
2 & \file{scripts/create\_project.tcl} has been dry-run but not yet run inside
    Vivado. First person to use it: confirm, then close this. & Joey \\
3 & Commit history contains a personal email address. It can only be cheaply
    replaced with the GitHub \texttt{noreply} address \emph{before} the first push;
    doing so changes every commit hash in this log. & Joey \\
\bottomrule
\end{tabularx}
\end{center}

% ============================================================================
\section{Lessons learned}
\begin{itemize}
  \item \textbf{``Programmed'' in Hardware Manager does not mean your design is on
        the chip}: it also shows for the factory demo loaded from flash. Use
        \emph{Program Device}. Each switch lighting its own LED means it's still the
        factory demo.
  \item \textbf{JTAG programming is volatile}: power-cycling restores the factory demo.
  \item \textbf{Vivado's simulator stops at 1\,\textmu s by default}; the project script
        sets it to run to completion (runtime \cmd{-all}).
  \item \textbf{Real data can hide bugs.} MNIST borders are blank, so pixel-skip
        bugs pass real images. Always add synthetic inputs that touch every path,
        and use mutation testing to prove the testbench can fail.
  \item \textbf{Weight layout is part of the interface.} Changing from neuron-major
        to pixel-major needed a new export file; the testbench checks one layout
        against the other.
\end{itemize}

% ============================================================================
\section{Entry template}\label{sec:template}
\begin{lstlisting}[language={[LaTeX]TeX}]
\subsection*{YYYY-MM-DD \quad Short title}
\begin{itemize}
  \item What changed (files, modules).
  \item How it was verified (testbench, checks, xsim or Icarus).
  \item Anything surprising or still open.
  \item Commit \texttt{abc1234}.
\end{itemize}
\end{lstlisting}

\end{document}
